\documentclass[11pt,a4paper]{article}

\usepackage[margin=1in]{geometry}
\usepackage{hyperref}
\usepackage{enumitem}
\usepackage{booktabs}
\usepackage{tabularx}
\usepackage{xcolor}
\usepackage{mdframed}

\hypersetup{colorlinks=true, linkcolor=blue!70!black, urlcolor=blue!70!black}

\title{ERC Indoor Team Start Here\\[0.3em]
\large High-level overview for everyone working on the project}
\author{Project Working Draft}
\date{\today}

\begin{document}

\maketitle
\begin{center}
\fbox{\parbox{0.9\textwidth}{
\textbf{HISTORICAL DOCUMENT} \\[4pt]
Early onboarding note from March 2026. Warning: several statements here are now
outdated. SuperGlue was never part of the final stack, and MBRA became the
primary indoor controller rather than a side experiment. \\[4pt]
\textit{For the current system, see: \texttt{CLAUDE.md},
\texttt{docs/INDEX.md}}
}}
\end{center}
\vspace{1em}

\section{What We Are Building}

We are building an indoor navigation system for the EarthRover Challenge in a
\textbf{known corridor environment}.

That means:

\begin{itemize}[nosep]
  \item the robot is not navigating an unknown building,
  \item we can collect data in the same corridor ahead of time,
  \item we can tune the system specifically for that corridor,
  \item our goal is a reliable competition system, not a general research demo.
\end{itemize}

\begin{mdframed}
\textbf{System choice:} we will use a \textbf{CosPlace + SuperGlue + graph}
pipeline for perception and planning, a \textbf{validated online local controller}
for short-horizon movement, and separate modules for \textbf{safety},
\textbf{checkpoint verification}, and \textbf{recovery}.

The MBRA / LogoNav repository remains an important reference:

\begin{itemize}[nosep]
  \item \textbf{MBRA} is primarily the short-horizon expert / relabeling model,
  \item \textbf{LogoNav} is the deployed-policy side of that project,
  \item neither should be assumed to be our indoor runtime controller without validation.
\end{itemize}
\end{mdframed}


\section{The Most Important File To Know}

\texttt{baseline.py} is the current baseline perception/planning implementation.

This file already contains the main ideas we are building on:

\begin{itemize}[nosep]
  \item CosPlace descriptors for place recognition,
  \item retrieval of matching corridor images,
  \item graph construction over reference images,
  \item action / connectivity graph generation,
  \item SuperPoint / SuperGlue geometric verification,
  \item debugging visualizations for graph and localization behavior.
\end{itemize}

\textbf{What \texttt{baseline.py} is:}

\begin{itemize}[nosep]
  \item the baseline for \textbf{perception and planning},
  \item the main file the perception and planning people should study first,
  \item a shared CLI script with two starting workflows:
  \item \texttt{build-db} for creating the reference database and graph artifacts,
  \item \texttt{query} for testing localization against that database.
\end{itemize}

\textbf{What \texttt{baseline.py} does not cover yet:}

\begin{itemize}[nosep]
  \item the full EarthRover runtime system,
  \item the final online controller for indoor edge-following,
  \item the safety layer,
  \item the sequential checkpoint mission loop.
\end{itemize}


\section{How The Full System Works}

At a high level, the robot loop is:

\begin{enumerate}[nosep]
  \item get the current front-camera image,
  \item localize that image in the corridor graph,
  \item plan which graph node to move toward next,
  \item give that nearby node image to the chosen online controller,
  \item let that controller propose a local motion command,
  \item run safety checks,
  \item send the safe command to the robot,
  \item verify whether the current checkpoint has been reached,
  \item if yes, switch to the next checkpoint.
\end{enumerate}


\section{Who Owns What}

\begin{tabularx}{\textwidth}{p{0.17\textwidth} p{0.28\textwidth} X}
\toprule
\textbf{Stack} & \textbf{Main job} & \textbf{Where to start} \\
\midrule
Perception & Figure out where the robot is in the corridor graph & Start with \texttt{baseline.py}; turn retrieval + verification into a live module \\
Planning & Decide which nearby graph node the robot should go to next & Start with graph construction and action-graph logic in \texttt{baseline.py} \\
Control & Move toward the chosen nearby node image & Start by deciding which online controller we will actually deploy, then build local command generation around it \\
Safety & Prevent unsafe commands from being sent & Build slow/stop/veto logic on top of control outputs \\
Checkpoint / Recovery & Decide when a checkpoint is reached and what to do when uncertain & Build goal confirmation and relocalization logic around the perception output \\
Integration & Connect all stacks into one runtime loop & Connect SDK, localization, planning, MBRA, safety, and mission state \\
\bottomrule
\end{tabularx}


\section{What Each Team Should Keep In Mind}

\subsection*{Perception}

Your job is to take the existing CosPlace + SuperGlue corridor pipeline and make
it work reliably in a live loop.

\subsection*{Planning}

Your job is to use the graph to choose the next nearby subgoal image.
Planning should keep global routing in the graph layer and hand off only nearby subgoals to the online controller.

\subsection*{Control}

Your job is to determine which controller we can actually run online for nearby
graph-node movement.

The MBRA / LogoNav codebase is the main reference here, but it should be read carefully:
\begin{itemize}[nosep]
  \item MBRA is the relabeling expert in the paper,
  \item LogoNav is the deployed-policy side,
  \item LogoNav in this repo is GPS/pose-conditioned, so it is not a drop-in indoor image-goal controller.
\end{itemize}

\subsection*{Safety}

Your job is to stop the robot from doing something dumb when the environment or
localization is uncertain.

\subsection*{Checkpoint / Recovery}

Your job is to prevent silent failure:
\begin{itemize}[nosep]
  \item prevent the system from claiming a checkpoint too early,
  \item prevent it from continuing confidently when it is lost.
\end{itemize}


\section{How To Build On This Baseline}

\begin{itemize}[nosep]
  \item Start from the existing baseline instead of rebuilding perception/planning from zero.
  \item Use the MBRA / LogoNav codebase as the main controller reference and adopt the online controller that proves strongest in real corridor testing.
  \item Treat monocular SLAM and other alternatives as possible enhancements if they prove stronger on real corridor data.
  \item Prioritize changes that improve reliability, recoverability, and safety on the robot.
  \item If someone has a better idea, the right standard is simple: show that it works better than the baseline on the real indoor route.
\end{itemize}


\section{Immediate Priority Order}

\begin{enumerate}[nosep]
  \item Understand \texttt{baseline.py}.
  \item Extract the reusable perception/planning pieces from it.
  \item Validate live localization on the corridor.
  \item Determine which online short-horizon controller is actually deployable.
  \item Validate that controller on nearby graph-node subgoals.
  \item Integrate planner $\rightarrow$ controller $\rightarrow$ safety.
  \item Add checkpoint verification and recovery.
  \item Then run full multi-checkpoint tests.
\end{enumerate}


\section{Bottom Line}

\begin{mdframed}
\begin{itemize}[nosep]
  \item \texttt{baseline.py} is the baseline perception/planning system.
  \item The online local controller still has to be selected and validated.
  \item Safety, checkpoint verification, and recovery are separate stacks.
  \item The main work now is integration and reliability on the real robot.
\end{itemize}
\end{mdframed}


\section{Related Documents}
\begin{itemize}
  \item \texttt{erc3\_full\_documentation.tex} --- single master guide for the complete project story and current architecture.
  \item \texttt{live\_indoor\_runtime\_story.tex} --- indoor evolution, MBRA integration, and checkpoint-step runtime behavior.
  \item \texttt{live\_outdoor\_ultra\_marathon\_story.tex} --- outdoor and marathon runtime evolution with safety-layer reasoning.
  \item \texttt{outdoor\_perception\_review.tex} --- depth/semantic perception findings and their runtime implications.
\end{itemize}

\end{document}
